\documentclass[conference]{IEEEtran}

\hyphenation{op-tical net-works semi-conduc-tor}

\begin{document}

\title{An AI-Assisted Framework for Automated Malware Reverse Engineering, Debugging, and Anti-Analysis Detection}

\author{
    \IEEEauthorblockN{Alexander Schoolcraft}
    \IEEEauthorblockA{
        Department of Computer Science and Cybersecurity\\
        College of Math and Sciences\\
        University of North Georgia\\
        Dahlonega, Georgia, 35033\\
        Email: adscho3199@ung.edu
    }
}

\maketitle

% ===========================================
% ABSTRACT
% ===========================================
\begin{abstract}
% Write this last. 150–250 words.
\end{abstract}

% ===========================================
% INTRODUCTION
% ===========================================
\section{Introduction}
Malware analysis remains a critical component of modern cybersecurity operations, providing the foundation for threat intelligence, incident response, and defensive tool development.

\subsection{Motivation and Problem Statement}
The modern malware landscape has grown increasingly sophisticated, leveraging advanced obfuscation, anti-debugging techniques, and complex execution patterns that challenge traditional analysis tools. As these threats evolve, understanding malicious binaries—through static analysis, dynamic instrumentation, and reverse engineering—remains labor-intensive and time-consuming. Security analysts routinely face significant backlogs of samples awaiting triage, where manual workflows and fragmented tools slow the investigation process and can delay detection, mitigation, and response.

This research aims to streamline and enhance the malware analysis workflow by developing an integrated system that combines established open-source tools—such as Ghidra, QEMU, and GDB—with modern Large Language Models (LLMs) and AI-driven reasoning components. The system is designed to unify static analysis, dynamic execution tracing, and automated debugging within a single coherent framework. By leveraging automation and AI-assisted interpretation, the proposed solution seeks to reduce analyst workload, accelerate the reverse engineering process, and generate detailed, structured reports describing malware behavior.

\subsection{Research Questions and Objectives}
This work is guided by the overarching goal of reducing the time, expertise, and manual effort required to analyze and understand malicious software. To achieve this, the following research questions were formulated:

\begin{itemize}
    \item \textbf{RQ1:} How can existing reverse-engineering and debugging tools (e.g., Ghidra, QEMU, GDB) be integrated into a unified analysis framework without sacrificing flexibility or depth of analysis?
    \item \textbf{RQ2:} To what extent can Large Language Models (LLMs) assist in reasoning about binary code, execution traces, and malware behavior when combined with static and dynamic analysis signals?
    \item \textbf{RQ3:} How can the system reliably detect and interpret anti-debugging and anti-analysis techniques, such as debug-register manipulation, self-relocating code, or control-flow disruptions?
    \item \textbf{RQ4:} What automation strategies can reduce analyst workload while still maintaining accuracy, interpretability, and reproducibility in malware analysis?
\end{itemize}

Based on these questions, the primary objectives of the research are as follows:

\begin{itemize}
    \item \textbf{O1:} Design and develop a modular architecture that orchestrates multiple open-source analysis tools within a single coherent system.
    \item \textbf{O2:} Implement LLM-driven reasoning modules capable of summarizing code behavior, identifying suspicious operations, and explaining malware functionality.
    \item \textbf{O3:} Incorporate mechanisms for detecting common and advanced anti-debugging techniques and integrating their signals into the analysis workflow.
    \item \textbf{O4:} Create an automated pipeline that executes malware samples in a controlled environment, collects relevant artifacts, and generates structured analysis reports.
    \item \textbf{O5:} Evaluate the system's effectiveness through case studies and controlled experiments that measure accuracy, completeness, and reduction in manual effort.
\end{itemize}

\subsection{Contributions}
This research aims to address the growing complexity and workload associated with malware analysis, reverse engineering, and debugging by integrating traditional program analysis tools with modern AI-driven reasoning. The anticipated contributions of this work are as follows:

\begin{itemize}
    \item \textbf{C1: Integrated Analysis Framework.}  
    The development of a unified architecture that orchestrates multiple open-source reverse-engineering and debugging tools—such as Ghidra, QEMU, and GDB—within a coherent, modular, and extensible system. This framework is designed to reduce fragmentation in current malware analysis workflows.

    \item \textbf{C2: LLM-Assisted Reasoning for Binary Analysis.}  
    The integration of Large Language Models (LLMs) to assist in interpreting static and dynamic artifacts, summarizing code behavior, identifying suspicious execution patterns, and supporting analysts with natural language explanations of complex operations.

    \item \textbf{C3: Automated Anti-Debugging and Anti-Analysis Detection.}  
    The inclusion of detection mechanisms for common and advanced anti-debugging techniques, such as debug-register manipulation, control-flow tampering, and self-relocating code, derived from recent research in attestation bypass and malware evasion.

    \item \textbf{C4: Automated Static and Dynamic Analysis Pipeline.}  
    The creation of an automated workflow capable of executing malware samples in a controlled environment, collecting relevant static, dynamic, and behavioral artifacts, and correlating them into a unified analysis record.

    \item \textbf{C5: Structured Behavioral Reporting.}  
    The generation of detailed, structured reports that combine program analysis signals with LLM-driven explanations, presenting malware behavior in a way that is accessible to analysts while preserving technical rigor.

    \item \textbf{C6: Evaluation Through Case Studies and Experiments.}  
    An empirical assessment of the system’s effectiveness, including case studies involving representative malware samples and controlled experiments measuring accuracy, interpretability, anti-debugging detection capability, and reduction in manual analysis effort.
\end{itemize}

\subsection{Paper Organization}
The remainder of this paper is structured as follows. Section~II provides background information on reverse engineering, debugging, and the tools and techniques commonly used in malware analysis. Section~III reviews related research and identifies the current gaps and limitations that motivate this work. Section~IV describes the design of the proposed system, beginning with a high-level architectural overview and followed by detailed descriptions of each core component. Section~V outlines the methodology used to integrate these components and explains how they collectively address the identified capability gaps. Section~VI describes the implementation of the system. Section~VII evaluates the system using representative analysis scenarios and discusses the resulting findings. Finally, Section~VIII summarizes the work, reflects on its implications, and identifies challenges and directions for future research.


% ===========================================
% BACKGROUND
% ===========================================
\section{Background}

This section provides the foundational technical context for the proposed system, including an overview of reverse engineering concepts, debugging mechanisms, anti-debugging techniques, and the open-source tools and AI models used in malware analysis.

\subsection{Reverse Engineering Fundamentals}

Reverse engineering (RE) is the process of analyzing compiled binaries to recover structural, semantic, and behavioral information about a program. Traditional RE workflows involve static analysis, which examines program structure without executing the binary, and dynamic analysis, which observes runtime behavior to identify execution paths, data manipulations, and system interactions. Analysts rely on control flow graphs (CFGs), data flow analysis, and function decompilation to infer program intent, especially in the absence of source code. Modern malware often introduces obfuscation, packing, or non-standard language runtimes that complicate these techniques.

\subsection{Debugging and Execution Monitoring}

Debugging is the act of observing and controlling a program during execution, typically by stepping through instructions, setting breakpoints, inspecting registers, or modifying memory. Tools such as GDB and LLDB interact with program state via well-defined debugging interfaces. Dynamic analysis may also be performed through whole-system emulators, such as QEMU, which allow analysts to instrument instructions or system calls in a controlled environment. The debugging architecture of modern processors, including x86 debug registers (DR0-DR7), hardware breakpoints, and single-step traps, plays a critical role in both legitimate debugging and adversarial anti-debugging.

\subsection{Anti-Debugging and Anti-Analysis Techniques}

Malware authors frequently employ anti-debugging techniques to hinder analysis and evade detection. These mechanisms include detecting hardware breakpoints, manipulating debug registers, abusing exception handlers, measuring timing discrepancies, self-modifying or self-relocating code, and tampering with control flow. Research such as self-relocating attacks based on debug-registers demonstrates how malware can erase or relocate itself during the attestation or debugging phases to avoid inspection. These behaviors pose significant challenges for traditional tools and motivate the need for integrated and automated detection mechanisms.

\subsection{Static and Dynamic Analysis Tools}

\subsubsection{Ghidra}
Ghidra is a widely used open-source reverse engineering framework that provides disassembly, decompilation, CFG reconstruction, and semantic analysis. Its extensibility and scripting capabilities allow analysts to automate repetitive tasks, extract program features, and perform structural analysis of malware.

\subsubsection{GDB}
The GNU Debugger (GDB) enables interactive debugging of binaries at runtime. It provides breakpoints, watchpoints, register inspection, memory analysis, and backtracing. GDB's machine interface (MI) and Python APIs make it suitable for integration into automated analysis pipelines.

\subsubsection{QEMU}
QEMU is a machine emulator capable of executing binaries in a virtualized environment while supporting binary translation, system-call logging, instrumentation, and hardware emulation. Its ability to observe low-level execution behavior makes it valuable for dynamic analysis and monitoring of malware that behaves differently under native debugging.

\subsection{Large Language Models and Program Understanding}

Large Language Models (LLMs) have recently shown promising capabilities in code summarization, error localization, symbolic reasoning, and debugging support. However, LLMs often struggle with deep semantic understanding, hallucination, or incomplete reasoning when operating without contextual signals. Hybrid approaches—combining LLMs with program analysis signals, execution traces, or probabilistic inference—have been shown to improve reliability and interpretability. These models serve as complementary components in malware analysis systems, providing natural-language explanations and high-level insights derived from raw program artifacts.

\subsection{Motivation for an Integrated System}

Traditional malware analysis relies on a collection of specialized but disconnected tools, requiring significant manual effort to orchestrate static analysis, dynamic monitoring, debugger control, anti-debugging detection, and report generation. As malware continues to adopt more sophisticated evasion and obfuscation strategies, analysts need systems that unify these capabilities and leverage automated reasoning. The background concepts presented in this section highlight both the power and limitations of existing tools, motivating the development of the integrated framework proposed in this work.

% ===========================================
% RELATED WORK
% ===========================================
\section{Related Work}

Research relevant to this work spans several areas, including AI-assisted debugging, reverse engineering and decompilation, anti-debugging and attestation bypass techniques, program comprehension using Large Language Models (LLMs), visualization systems for binary analysis, and educational or human-centered debugging support. This section summarizes key developments in these domains and highlights the gaps that motivate the proposed system.

% -------------------------------------------
\subsection{LLM-Assisted Debugging and Code Reasoning}

Recent work has explored the use of LLMs to assist in program debugging, error localization, and code explanation. ChatDBG~\cite{levin2024chatdbg} integrates LLMs with existing debuggers to support natural-language queries and automated stepping, demonstrating the potential of AI to augment traditional debugging workflows. AutoSD~\cite{kang2023autosd} applies LLMs to automated scientific debugging, combining hypothesis generation with program instrumentation to explain errors and propose fixes. These systems illustrate that LLMs can provide useful insights, but they largely operate on source code and do not integrate static, dynamic, or binary-level signals. Furthermore, they do not address anti-analysis or malware-specific challenges.

% -------------------------------------------
\subsection{LLMs for Reverse Engineering and Program Comprehension}

Several works investigate the role of LLMs in reverse engineering tasks. The ProCURE framework~\cite{ren2025procure} enhances conceptual reasoning in LLMs through counterfactual code generation, improving understanding of control flow and data dependencies. Beyond C/C++~\cite{zhang2025beyondccpp} highlights the limitations of current tools when analyzing binaries produced from modern languages such as Rust and Go, and proposes hybrid probabilistic–LLM techniques to improve semantic inference. Visualization agents such as those presented in the CogBRE VR system~\cite{brown2025llmvisre} demonstrate how LLMs can drive immersive analysis environments, though these works focus primarily on educational or interactive analysis rather than automated malware workflows.

Other studies, including LLM-based software security analyses~\cite{vaddiparthy2025llmsecurity}, explore how AI models can support vulnerability and malware analysis, but they often lack integration with real debugging or emulation tools.

% -------------------------------------------
\subsection{Reverse Engineering Tools and Decompiler Benchmarks}

Tools such as Ghidra, GDB, QEMU, angr, and capa form the backbone of modern malware analysis. DecompileBench~\cite{gao2025decompilebench} evaluates commercial and open-source decompilers at scale, demonstrating strengths and weaknesses in decompilation accuracy and understandability. While these tools provide rich program analysis capabilities, they require significant manual coordination, and no existing framework combines their outputs with LLM-driven reasoning in a coherent, automated system.

% -------------------------------------------
\subsection{Anti-Debugging, Attestation, and Evasion Techniques}

Malware authors frequently deploy anti-debugging mechanisms to frustrate analysis. Debug-register-based self-relocating attacks (DRSA)~\cite{bypassremoteattestation2025} illustrate how adversaries can manipulate debug exceptions to evade attestation and integrity checks. Measurements of anti-debugging techniques across Windows and Linux~\cite{antidebugmeasure2025} reveal the diversity and sophistication of modern evasion behaviors. Branch-monitoring techniques for security~\cite{botacin2018branchmonitoring} demonstrate hardware-assisted detection strategies but require specialized monitoring and do not directly integrate into automated debugging pipelines. These works highlight the persistent challenge of detecting and interpreting anti-analysis behavior in real malware.

% -------------------------------------------
\subsection{Visualization and Educational Support for Debugging}

Immersive visualization environments and instructional systems also inform this research. VR-based visualization of binary analysis~\cite{brown2025llmvisre} enables novel forms of structural comprehension, while LLM-driven educational tools~\cite{k12llmdebugging2025,usemewisely2025} show that AI can enhance debugging skills, prompt formulation, and conceptual understanding. Although valuable, these systems focus on pedagogy or visualization rather than automated malware analysis.

% -------------------------------------------
\subsection{Summary and Research Gap}

Existing work demonstrates significant advances in reverse engineering, debugging automation, anti-evasion detection, and LLM-assisted reasoning. However, no current system integrates:

\begin{itemize}
    \item multi-tool analysis workflows (e.g., Ghidra + QEMU + GDB),
    \item LLM-driven program interpretation at the binary level,
    \item automated anti-debugging and anti-analysis detection,
    \item correlated static and dynamic analysis pipelines, and
    \item structured malware behavior reporting.
\end{itemize}

These gaps motivate the integrated framework proposed in this research, which aims to unify these elements into a single, coherent, AI-assisted malware analysis system.


% ===========================================
% SYSTEM DESIGN
% ===========================================
\section{System Design}

\subsection{Architecture Overview}
% High-level description + reference a figure.

\subsection{Core Components}
\subsubsection{Static Analysis Module}
\subsubsection{Dynamic Execution Monitor}
\subsubsection{Anti-Debug Detection Engine}
\subsubsection{LLM Reasoning and Tool Orchestration}
\subsubsection{Probabilistic Inference Layer}
\subsubsection{User Interface / Visualization Layer}

\subsection{Data Flow}
% Pipeline description.

\subsection{Design Assumptions and Threat Model}
% Based on anti-debug/attestation papers.

% ===========================================
% METHODOLOGY
% ===========================================
\section{Methodology}

\subsection{LLM Integration Strategy}
% Prompting, safety, hallucination mitigation.

\subsection{Debugger Integration}
% GDB/LLDB APIs, trace extraction, etc.

\subsection{Static/Dynamic Analysis Techniques}
% How you collect signals for reasoning.

\subsection{Anti-Debugging Detection Methods}
% Derived from DRSA, branch monitoring, etc.

% ===========================================
% IMPLEMENTATION
% ===========================================
\section{Implementation}

\subsection{Technologies and Frameworks Used}
% Python/C++ hybrid, MCP servers, OpenAI API, etc.

\subsection{System Modules and Code Structure}
% Folder layout, module responsibilities.

\subsection{Integration with External Tools}
% Ghidra, angr, capa, CUDA (if used), etc.

\subsection{Limitations of Implementation}
% Practical constraints.

% ===========================================
% EXPERIMENTAL EVALUATION
% ===========================================
\section{Experimental Evaluation}

\subsection{Experimental Setup}
% Hardware, software, LLMs, tools.

\subsection{Test Programs and Benchmarks}
% DecompileBench, synthetic binaries, anti-debug samples.

\subsection{Evaluation Metrics}
% Accuracy, reasoning quality, detection rate, latency.

\subsection{Case Studies}
% Walk through specific debugging tasks.

\subsection{Results and Analysis}
% Numbers, graphs, findings.

% ===========================================
% DISCUSSION
% ===========================================
\section{Discussion}

\subsection{Interpretation of Results}
% What do they mean?

\subsection{Implications for Debugging and RE}
% Big-picture takeaways.

\subsection{Threats to Validity}
% Internal, external, construct validity, etc.

% ===========================================
% CHALLENGES AND FUTURE WORK
% ===========================================
\section{Challenges and Future Work}

\subsection{Technical Challenges Encountered}
% Weak tool documentation, API weirdness, data noise.

\subsection{Potential Improvements}
% Better LLM tuning, larger datasets, etc.

\subsection{Long-Term Research Directions}
% Fully autonomous debugger, formal reasoning, etc.

% ===========================================
% CONCLUSION
% ===========================================
\section{Conclusion}
% Final summary + how your system advances the field.

% ===========================================
% ACKNOWLEDGMENT
% ===========================================
\section*{Acknowledgment}
% Optional.

% ===========================================
% REFERENCES
% ===========================================
%\nocite{*}
\bibliographystyle{IEEEtran}
\bibliography{cites}

\end{document}
